\begin{table}[!htbp]
\centering
\small
\renewcommand{\arraystretch}{1.2}
\begin{tabular}{p{0.10\linewidth} p{0.48\linewidth} p{0.34\linewidth}}
\hline
\textbf{Label} &
\textbf{Assumption} &
\textbf{Role in the comparison}
\\
\hline

A1 &
The physical circuit noise is described by local Markovian CPTP channels
acting at prescribed circuit locations. &
Defines the physical noise model used in the resource analysis. The
common-reference-class result additionally specializes this setting to the
Pauli-noise and circuit conditions of
Appendix~\ref{app:pauli_noise_complexity}.
\\

A2 &
For the realized training sample $S$, fixed margin threshold $\mu$, and
selected PQC architecture, there exists
$h_\theta\in\mathcal H_d^{\mathrm{noiseless}}$
such that
$\widehat R_{S,\mu}(h_\theta)\leq\varepsilon_{\mathrm{mar}}$. &
Ensures that the common empirical margin-loss target is attainable by the
ideal reference class.
\\

A3 &
The ideal training procedure reaches the target
$\widehat R_{S,\mu}\leq\varepsilon_{\mathrm{mar}}$
within a finite expectation-value evaluation budget
$N_{\mathrm{noiseless}}$
with probability at least $1-\delta_{\mathrm{opt}}$. &
Defines the ideal reference resource budget against which PEC and boosting
are compared.
\\

\hline
\end{tabular}
\caption{Assumptions common to the ideal reference problem and the
PEC--boosting comparison.}
\label{tab:common_assumptions}
\end{table}
